\documentclass[12pt,a4paper]{article}

\usepackage[margin=1in]{geometry}
\usepackage{newtxtext,newtxmath}
\usepackage[hidelinks]{hyperref}
\usepackage{setspace}
\usepackage[ngerman]{babel}
\usepackage{changepage}

\setlength{\parskip}{0.5\baselineskip}
\setlength{\parindent}{0pt}

\pagestyle{empty}
\onehalfspacing

\begin{document}

\textbf{Quin Barnard}
\hfill
\textbf{Mittwoch, den 26. August 2026}\\
\textbf{\#27902277}\\
\textbf{Deutsch 348: Abschlussarbeit}\\
\textbf{Herr Professor Ingo Müller}\\
\textbf{Essay-Skizze:} Kafkas \textit{Betrachtung}\\

Bei den Erzählungen von Kafka sind die Affekte Angst, Scham und Schuld am wichtigsten; trotzdem sind Unbehagen, Unsicherheit, Verlegenheit und Unentschlossenheit hervorstechend. Da diese stark mit denjenigen verbunden sind, könnte sich eine Analyse die sich auf die weniger bemerkten Affekte schränkt lohnen.\par

Im Band \textit{Betrachtung} handeln die Hauptpersone kaum ohne von Unsicherheit ergriffen zu werden. Der Protagonist bei \textit{Kinder auf der Landstraße} wartet auf die Kinder und seufzt wie von einem „nie gut zu [machenden] Unglück“ berührt als diejenige erscheinen. Es scheint fast als wäre der Protagonist allein nicht zu seinem eigenen Glücklichsein fähig, weil er aufgemuntert werden muss. Als die Kinder eine Strecke gelaufen sind, teilt sich der Gruppe in zwei, was den Protagonist wieder beunruhigt und Konflikt verursacht. Der Leser bekommt hauptsächlich die Gefühlen und Gedanken des Protagonisten mit und nicht die Tatsachen als deren Ursachen. Da es die Rede von einem Ich-Erzähler ist, ist es dem Leser nicht möglich zu schätzen welche Ereignisse inwiefern real sind. Der Leser weiß nur, dass der Protagonist unbequem ist, und in gewissermaßen seine eigene Identität bezweifelt als sich die Mitglieder der Gruppe streiten.\par

Ähnlich verhält es sich bei \textit{Unglücklichsein}. Der Protagonist versucht das Gespräch von einem ernsthaften Vorwurf abzulenken indem er sagt: 

\begin{adjustwidth}{1.5cm}{1.5cm}
    \begin{spacing}{1}
        \begin{center}
            \small
                \textit{Wie? Ich hätte Ihnen gedroht? Aber ich bitte Sie. Ich bin ja so froh, daß Sie endlich hier sind. Ich sage ›endlich‹ weil es schon so spät ist. Da ist es möglich, daß ich in der Freude so durcheinander gesprochen habe und daß Sie es gerade so verstanden haben.}
        \end{center}
    \end{spacing}
\end{adjustwidth}

Es kommt dem Leser vor, als ob die Hauptperson etwa Verlegen wäre und daher das Kind zu beschwichtigen versucht. Die Verlegenheit stammt daraus, dass die Hauptperson und das Kind verschiedene Meinungen haben über was geschehen und gesagen ist. Der Leser weiß wieder nichts als was der Erzähler sagt, und kann daher nicht wissen welche Partei Recht hat.\par

Sowohl die Epistemologie als auch die Ontologie sind nützlich beim Deuten dieser beiden Erzählungen. Die Frage, wie und in wiefern das Subjekt die „Realität“ anders empfindet wenn es unsicher, verlegen, unentschlossen usw. ist, ist kaum ohne sie möglich zu beantworten. Die Theorie von Roland Barthes wie von Ulrike Vedder\footnote{Vedder, Ulrike: Das Rätsel der Objekte: Zur literarischen Epistemologie von Dingen. Eine Einführung. In: Zeitschrift für Germanistik 22 (2012), H. 1, S.7-16} erwähnt könnte helfen herauszufinden, wie und sogar ob man von einer externen Realität sprechen kann, wenn man sie nie unmittelbar herangehen kann. Zu Kafkas Werke gibt es viele Sekondärtexte die die Analyse erleichtern können, wie \textit{Uncertainty, Realism, and the Self in Kafka}\footnote{Gosetti-Ferencei, Jennifer: Uncertainty, Realism and the Self in Kafka. In: Duttlinger, Carolin; Hilliard, K.F. und Louth, Charlie (Hg.): From the Enlightenment to Modernism: Three Centuries of German Literature. Oxford 2021, ohne Seiten Angabe} und \textit{The Dynamics of Narration in Betrachtung, “Das Urteil,” and Kafka's Reflections on Writing.}\footnote{Speirs, Ronald: The Dynamics of Narration in Betrachtung, “Das Urteil,” and Kafka's Reflections on Writing. In: Lothe, Jakob; Sandberg, Beatrice und Speirs, Ronald (Hg.): Theory and Interpretation of Narrative. Columbus 2011, S.196-232}

\end{document}

